\appendix
\section{שחזור והרצה}

הדוח והמחברת קוראים תוצאות machine-readable מה־artifacts המאומתים, ולכן ניתן לעיין בהם בלי לאמן מחדש את כל המודלים. הריצות המרכזיות הן:

\begin{itemize}
    \item \code{experiments_extended/extended_20260811_121957/}
    \item \code{experiments_supervised/supervised_20260811_133344/}
\end{itemize}

סביבת הריצה המאומתת הייתה Python 3.11.15. בריצת ה־HMM נעשה שימוש ב־hmmlearn 0.3.3, NumPy 2.4.6, pandas 3.0.3 ו־scikit-learn 1.9.0.

להרצה מחדש:

\begin{english}
\begin{verbatim}
python -m pip install -r requirements.txt
python run_extended_analysis.py
python analyze_extended_context.py \
  --run-dir experiments_extended/<run_id>
python run_supervised_baselines.py
jupyter nbconvert --to notebook --execute \
  HMM_Market_Regimes_Project.ipynb \
  --output /tmp/HMM_verified.ipynb \
  --ExecutePreprocessor.timeout=600
\end{verbatim}
\end{english}

כל ריצה חדשה נכתבת לתיקייה timestamped חדשה ואינה אמורה לדרוס את התוצרים שעליהם מבוסס הדוח. הדבר חשוב במיוחד משום שמקור הנתונים חיצוני ונתוני Yahoo Finance יכולים להתעדכן לאורך זמן.
